% Vorlage für ein Aufgabenblatt (schwarz-weiß, druckfertig)
% Kompilieren: pdflatex -interaction=nonstopmode template.tex
\documentclass[11pt,a4paper]{article}

\usepackage[T1]{fontenc}
\usepackage[utf8]{inputenc}
\usepackage[ngerman]{babel}
\usepackage[left=2cm,right=2cm,top=2cm,bottom=2.5cm]{geometry}
\usepackage{amsmath,amssymb}
\usepackage{tikz}
\usetikzlibrary{arrows.meta}
\usepackage{array,booktabs}
\usepackage{enumitem}
\usepackage{needspace}
\usepackage{fancyhdr}
\renewcommand{\familydefault}{\sfdefault}

% --- Fußzeile: Fach, Jahrgang, Dokumenttyp, Dokument-ID ---
\pagestyle{fancy}
\fancyhf{}
\renewcommand{\headrulewidth}{0pt}
\renewcommand{\footrulewidth}{0.4pt}
\fancyfoot[L]{\footnotesize Mathematik $\cdot$ Klasse 8 $\cdot$ Übungsblatt}
\fancyfoot[C]{\footnotesize $\cdot$ 0405 $\cdot$}
\fancyfoot[R]{\footnotesize \thepage}

% --- Antwortzeilen ---
\newcommand{\AL}{\noindent\rule{\linewidth}{0.3pt}\par\vspace{5mm}}
\newcommand{\antwortbereich}[1]{\vspace{3mm}\ifcase#1\or\AL\or\AL\AL\or\AL\AL\AL%
\or\AL\AL\AL\AL\or\AL\AL\AL\AL\AL\or\AL\AL\AL\AL\AL\AL\or\AL\AL\AL\AL\AL\AL\AL%
\or\AL\AL\AL\AL\AL\AL\AL\AL\or\AL\AL\AL\AL\AL\AL\AL\AL\AL\fi\vspace{1mm}}

% --- Abschnittsheader ohne Farbe ---
\newcommand{\abschnitt}[1]{\Needspace{4\baselineskip}\vspace{4mm}%
  \noindent\textbf{#1}\par\vspace{1mm}\noindent\rule{\linewidth}{0.5pt}\par\vspace{2mm}}

% --- Aufgabenumgebung: hält Teilaufgaben zusammen ---
\newcounter{aufg}
\newcommand{\aufgabe}[1]{\Needspace{5\baselineskip}\vspace{4mm}%
  \stepcounter{aufg}\noindent\textbf{Aufgabe \theaufg}\quad #1\par\vspace{1mm}}

\begin{document}

% --- Kopf: NUR der Thementitel ---
\begin{center}
  {\Large\bfseries Lineare Funktionen: Nullstellen und Schnittpunkte}
\end{center}
\vspace{4mm}

\aufgabe{Berechne die Nullstelle der Funktion.}
\begin{enumerate}[label=\alph*), itemsep=6mm]
  \item $f(x) = 2x - 6$
  \item $f(x) = -3x + 12$
  \item $f(x) = 0{,}5x - 2{,}5$
\end{enumerate}

\aufgabe{Lies die Steigung und den y-Achsenabschnitt aus dem Schaubild ab und stelle die Funktionsgleichung auf.}
\begin{center}
\begin{tikzpicture}[scale=0.65]
  \draw[very thin,gray!30] (-4.5,-4.5) grid (5.5,5.5);
  \draw[-{Stealth[length=6pt]},thick] (-4.8,0) -- (5.8,0) node[right,font=\small] {$x$};
  \draw[-{Stealth[length=6pt]},thick] (0,-4.8) -- (0,5.8) node[above,font=\small] {$y$};
  \foreach \x in {-4,...,5} {
    \draw[thin] (\x,0.08) -- (\x,-0.08);
    \node[below,font=\scriptsize] at (\x,-0.15) {\x};
  }
  \foreach \y in {-4,...,5} {
    \draw[thin] (0.08,\y) -- (-0.08,\y);
    \node[left,font=\scriptsize] at (-0.15,\y) {\y};
  }
  \draw[very thick] (-2.5,-4) -- (4.5,3);
\end{tikzpicture}
\end{center}
\antwortbereich{2}

\aufgabe{Fülle die Tabelle aus.}
\begin{center}
\begin{tabular}{|c|c|c|c|}
\hline
$x$ & $-2$ & $0$ & $3$ \\ \hline
$f(x) = 3x - 1$ & & & \\[4mm] \hline
\end{tabular}
\end{center}

\aufgabe{Bestimme den Schnittpunkt der beiden Geraden $g: y = 2x - 1$ und $h: y = -x + 5$.}
\antwortbereich{4}

\aufgabe{Zeige rechnerisch, dass die Geraden $g: y = 4x + 2$ und $h: y = 4x - 3$ keinen Schnittpunkt besitzen. Begründe deine Rechnung mit einem Satz.}
\antwortbereich{4}

\newpage

\abschnitt{Lösungen}

\noindent\textbf{Aufgabe 1}\par
a) $2x - 6 = 0 \Rightarrow 2x = 6 \Rightarrow x = 3$\par
b) $-3x + 12 = 0 \Rightarrow -3x = -12 \Rightarrow x = 4$\par
c) $0{,}5x - 2{,}5 = 0 \Rightarrow 0{,}5x = 2{,}5 \Rightarrow x = 5$\par

\vspace{3mm}
\noindent\textbf{Aufgabe 2}\par
Die Gerade verläuft durch $(-2{,}5\,|\,-4)$ und $(4{,}5\,|\,3)$, also
$m = \dfrac{3 - (-4)}{4{,}5 - (-2{,}5)} = \dfrac{7}{7} = 1$; der y-Achsenabschnitt ist $b = -1{,}5$.
Damit gilt $f(x) = x - 1{,}5$.

\vspace{3mm}
\noindent\textbf{Aufgabe 3}\par
$f(-2) = -7$, \quad $f(0) = -1$, \quad $f(3) = 8$

\vspace{3mm}
\noindent\textbf{Aufgabe 4}\par
$2x - 1 = -x + 5 \Rightarrow 3x = 6 \Rightarrow x = 2$, \quad $y = 2 \cdot 2 - 1 = 3$.
Schnittpunkt: $S(2\,|\,3)$.

\vspace{3mm}
\noindent\textbf{Aufgabe 5}\par
$4x + 2 = 4x - 3 \Rightarrow 2 = -3$. Das ist eine falsche Aussage, die Gleichung hat also keine Lösung.
Beide Geraden haben die Steigung $m = 4$, sie sind parallel und schneiden sich daher nicht.

\end{document}
